\chapter{Conclusiones}
El desarrollo del Sistema de Gestión de Turnos Académicos ha permitido construir una solución tecnológica integral que responde a las necesidades específicas de la Universidad Nacional de Salta y, en particular, de su Observatorio. A lo largo de este proyecto, se ha logrado digitalizar un proceso que anteriormente dependía de métodos manuales o fragmentados, sentando las bases para una gestión institucional más moderna y transparente.

\section{Logros alcanzados}
\begin{itemize}
    \item \textbf{Centralización y Orden:} Se logró implementar un motor de reservas capaz de gestionar la jerarquía académica de la UNSa, permitiendo que cada dependencia tenga autonomía sobre sus trámites y capacidades de atención.
    \item \textbf{Optimización del Observatorio:} La funcionalidad de turnos grupales y el control dinámico de cupos aseguran que las visitas institucionales se realicen de manera organizada, evitando aglomeraciones y mejorando la calidad de la divulgación científica.
    \item \textbf{Seguridad y Confianza:} El uso de estándares de industria (RBAC, CSRF, Hashing) garantiza la protección de los datos de los solicitantes y la integridad de la administración interna.
\end{itemize}

\section{Trabajos Futuros}
A pesar de la robustez de la solución actual, se identifican líneas de mejora que podrían expandir el impacto del sistema:
\begin{itemize}
    \item \textbf{Integración con SIU-Guaraní:} Vincular el sistema con el ecosistema de gestión de alumnos de la universidad para precargar datos de estudiantes de forma automática.
    \item \textbf{Aplicación Móvil:} Desarrollar una aplicación nativa o PWA que permita a los usuarios gestionar sus turnos y recibir notificaciones \textit{push} en tiempo real.
    \item \textbf{Módulo de Encuestas:} Incorporar un sistema de feedback post-atención para medir la satisfacción de los visitantes del Observatorio.
\end{itemize}

Finalmente, la utilización del stack tecnológico TALL (en parte) y Laravel 8 ha demostrado ser una elección acertada, permitiendo un desarrollo ágil y escalable que cumple con los requerimientos técnicos y funcionales planteados al inicio de este seminario técnico profesional.
